%% LyX 2.5.1 created this file.  For more info, see https://www.lyx.org/.
%% Do not edit unless you really know what you are doing.
\documentclass[12pt,a4paper,english]{article}
\usepackage[T1]{fontenc}
\usepackage[utf8]{inputenc}
\setlength{\parindent}{0bp}
\usepackage{color}

\makeatletter

%%%%%%%%%%%%%%%%%%%%%%%%%%%%%% LyX specific LaTeX commands.
\special{papersize=\the\paperwidth,\the\paperheight}


%%%%%%%%%%%%%%%%%%%%%%%%%%%%%% Textclass specific LaTeX commands.
\newenvironment{lyxcode}
	{\par\begin{list}{}{
		\setlength{\rightmargin}{\leftmargin}
		\setlength{\listparindent}{0pt}% needed for AMS classes
		\raggedright
		\setlength{\itemsep}{0pt}
		\setlength{\parsep}{0pt}
		\normalfont\ttfamily}%
	 \item[]}
	{\end{list}}

%%%%%%%%%%%%%%%%%%%%%%%%%%%%%% User specified LaTeX commands.
\usepackage{xcolor}

\makeatother

\usepackage{babel}
\begin{document}
\title{Konvertor trofaznih tokova snaga: polarne koordinate}
\author{SREES\_2026\_Bajramovic\_TPFPC}
\date{Sarajevo, 2026.}

\maketitle
Univerzitet u Sarajevu

Elektrotehnički fakultet

Odsjek za automatiku i elektroniku

Predmet: SREES

Akronim: TPFPC, Three Phase Flow Polar Converter

Student: Bajramović Anesa
\begin{abstract}
Ovaj izvještaj opisuje projekat SREES\_2026\_Bajramovic\_TPFPC, C++/natID
aplikaciju za konverziju trofaznih kompleksnih snaga iz pravougaonog
oblika P+jQ u polarni oblik |S| i ugao izražen u stepenima. Projekat
obuhvata glavni natID GUI, balanced i unbalanced režim rada, tok podataka
kroz dense i sparse matrice, worker/progress threadove, obični DLL
plugin i dTwin import plugin koji generiše .dmodl model. Projekat
nije puni load-flow solver; fokus je isključivo na konverziji poznatih
P/Q vrijednosti u polarne koordinate.
\end{abstract}

\section{Uvod}

Kompleksna snaga se u elektroenergetskim sistemima često zapisuje
u obliku P+jQ, gdje P predstavlja aktivnu, a Q reaktivnu komponentu.
Za prikaz, poređenje i dalju obradu rezultata često je praktičnije
koristiti polarni oblik, odnosno magnitudu i ugao. U trofaznim sistemima
dodatno je važno podržati i balansirane i nebalansirane ulaze, jer
realni sistemi ne moraju imati jednake vrijednosti po fazama.

TPFPC je zbog toga zamišljen kao čisti konverter: korisnik unosi P/Q
vrijednosti po fazama, a aplikacija računa magnitude, uglove i totalne
vrijednosti. U projektu se ne implementira Newton-Raphson solver i
ne koristi se Ybus demo kao dio finalne funkcionalnosti.

{\large\textcolor{red}{{[}OVDJE UBACITI SLIKU: GUI glavne natID aplikacije{]}}}{\large\par}

\section{Cilj projekta}

Cilj projekta je izrada C++/natID alata za trofaznu konverziju P+jQ
u polarne koordinate. Projekat obuhvata:
\begin{itemize}
\item C++ implementaciju matematičke logike za konverziju.
\item natID GUI za unos i prikaz rezultata.
\item Balanced i unbalanced trofazni unos.
\item Konverziju u posebnom worker threadu.
\item Progress indikator koji se osvježava iz posebnog progress threada.
\item natID dense i sparse matrice kao stvarni dio toka podataka.
\item Obični DLL plugin za izdvajanje core konverzije.
\item dTwin import plugin koji generiše .dmodl model.
\item CMake build konfiguraciju.
\end{itemize}

\section{Teorijska osnova}

Za svaku fazu se kompleksna snaga zapisuje kao:
\begin{lyxcode}
Sa~=~Pa~+~jQa

Sb~=~Pb~+~jQb

Sc~=~Pc~+~jQc
\end{lyxcode}
Konverzija jedne faze iz pravougaonog u polarni oblik računa se formulama:
\begin{lyxcode}
|S|~=~sqrt(P\textasciicircum 2~+~Q\textasciicircum 2)

angleDeg~=~atan2(Q,~P)~{*}~180~/~pi
\end{lyxcode}
Za totalne vrijednosti koristi se zbir po fazama:
\begin{lyxcode}
Ptotal~=~Pa~+~Pb~+~Pc

Qtotal~=~Qa~+~Qb~+~Qc

|Stotal|~=~sqrt(Ptotal\textasciicircum 2~+~Qtotal\textasciicircum 2)

angleTotal~=~atan2(Qtotal,~Ptotal)~{*}~180~/~pi
\end{lyxcode}

\section{Balanced i unbalanced režim}

U balanced režimu korisnik unosi samo vrijednosti za fazu A. Vrijednosti
faza B i C se zatim interno popunjavaju istim P/Q vrijednostima. Ovaj
režim je koristan kada su sve tri faze jednake i kada nije potrebno
ručno ponavljati isti unos.

U unbalanced režimu korisnik unosi P/Q vrijednosti za sve tri faze
posebno. Ovaj režim odgovara realnijim situacijama u kojima faze A,
B i C mogu imati različite aktivne i reaktivne komponente.

{\large\textcolor{red}{{[}OVDJE UBACITI SLIKU: Balanced režim{]}}}{\large\par}

{\large\textcolor{red}{{[}OVDJE UBACITI SLIKU: Unbalanced režim{]}}}{\large\par}

\section{Arhitektura aplikacije}

Tok podataka u aplikaciji je organizovan tako da se GUI unos prvo
pretvara u matricni model, zatim se provodi konverzija i na kraju
se rezultati prikazuju u GUI-ju ili zapisuju u dTwin .dmodl izlaz.
\begin{lyxcode}
GUI~input

->~dense~input~matrix~3x2

->~sparse~input~matrix~3x2

->~PolarConverter

->~dense~output~matrix~3x2

->~sparse~output~matrix~3x2

->~GUI~rezultat~/~dTwin~.dmodl~output
\end{lyxcode}
{\large\textcolor{red}{{[}OVDJE UBACITI DIJAGRAM: tok podataka GUI
-> matrice -> konverter -> rezultat{]}}}{\large\par}

\section{Implementacija u C++/natID}

Implementacija je podijeljena u više manjih fajlova, tako da je matematička
logika odvojena od GUI-ja, matrica i plugin slojeva.
\begin{description}
\item [{PowerFlowTypes.h}] Definiše osnovne strukture za pravougaoni i
polarni zapis snage.
\item [{PolarConverter.h/.cpp}] Implementira formulu za magnitudu, ugao
u radijanima i ugao u stepenima.
\item [{PowerFlowMatrixModel.h/.cpp}] Povezuje ulazne P/Q vrijednosti,
dense matrice, sparse matrice, konverziju i totalne rezultate.
\item [{PowerFlowView.h/.cpp}] Implementira glavni natID GUI, validaciju
unosa, balanced/unbalanced režim, worker thread i progress thread.
\item [{MatricesView.h/.cpp}] Predviđen je za prikaz matrica u posebnom
tabu.
\item [{TpfpcPluginApi.h/.cpp}] Definiše obični DLL C API sloj za pozivanje
core konverzije.
\item [{TpfpcDTwinPlugin.h/.cpp}] Implementira dTwin import plugin, prozor
za unos i generisanje .dmodl modela.
\end{description}

\subsection{Kodni isječak: formula u PolarConverter}

Fajl: PolarConverter.cpp
\begin{lyxcode}
polarPower.magnitude~=~std::sqrt(power.p~{*}~power.p~+~power.q~{*}~power.q);

polarPower.angleRad~=~std::atan2(power.q,~power.p);

polarPower.angleDeg~=~polarPower.angleRad~{*}~180.0~/~pi;
\end{lyxcode}
Ovaj kod direktno implementira osnovnu matematičku formulu projekta.
Bitan je zato što se na njemu zasnivaju rezultati za sve faze, totalne
vrijednosti, plugin rezultat i dTwin .dmodl izlaz.

\subsection{Kodni isječak: balanced/unbalanced logika}

Fajl: PowerFlowView.cpp
\begin{lyxcode}
if~(isBalancedMode())

\{

~~~~input.phaseB~=~input.phaseA;

~~~~input.phaseC~=~input.phaseA;

~~~~return~true;

\}
\end{lyxcode}
Ovaj isječak pokazuje da se u balanced režimu vrijednosti faze A koriste
i za faze B i C. Ako balanced režim nije uključen, kod nastavlja čitati
ručne P/Q vrijednosti za faze B i C, što pokriva unbalanced režim.

\section{Dense i sparse matrice}

Dense input matrix je matrica dimenzija 3x2. Redovi predstavljaju
faze A, B i C, a kolone predstavljaju P i Q. Dense output matrix je
također 3x2, ali kolone predstavljaju magnitude i angleDeg.

Sparse input i sparse output matrice čuvaju nenulte vrijednosti iz
ulaznog i izlaznog toka podataka. Ove matrice nisu demonstrativni
Ybus i ne predstavljaju load-flow mrežni model. One su dio stvarnog
toka konverzije P/Q vrijednosti u polarne koordinate.

{\large\textcolor{red}{{[}OVDJE UBACITI SLIKU: Matrices tab{]}}}{\large\par}

\subsection{Kodni isječak: dense input/output i sparse nnz}

Fajl: PowerFlowMatrixModel.cpp
\begin{lyxcode}
auto~input~=~result.inputMatrix.getManipulator();

input(kPhaseA,~kPColumn)~=~rectangularPower.phaseA.p;

input(kPhaseA,~kQColumn)~=~rectangularPower.phaseA.q;

const~auto~phaseA~=~\_converter.convertPhase(\{input(kPhaseA,~kPColumn),~input(kPhaseA,~kQColumn)\});

writePolarToMatrix(result.polarMatrix,~kPhaseA,~phaseA);

addSparseValue(sparseInput,~kPhaseA,~kPColumn,~input(kPhaseA,~kPColumn),~result.sparseInputNonZeroCount);

addSparseValue(sparseOutput,~kPhaseA,~kMagnitudeColumn,~polar(kPhaseA,~kMagnitudeColumn),~result.sparseOutputNonZeroCount);
\end{lyxcode}
Ovaj kod pokazuje kako se ulaz prvo zapisuje u dense matricu, zatim
se konvertuje preko PolarConverter-a i rezultat upisuje u dense output
matricu. Pozivi addSparseValue broje i zapisuju samo nenulte vrijednosti,
što je bitno za validaciju sparse toka podataka.

\section{Threading i progress indikator}

Glavna aplikacija koristi odvojeni worker thread za konverziju i progress
thread za osvježavanje indikatora napretka. Pošto se GUI elementi
ne smiju direktno mijenjati iz pozadinskih threadova, update se šalje
u glavni GUI thread preko gui::thread::asyncExecInMainThread.

{\large\textcolor{red}{{[}OVDJE UBACITI DIJAGRAM: worker thread +
progress thread + GUI main thread{]}}}{\large\par}

\subsection{Kodni isječak: worker/progress thread i thread-safe GUI update}

Fajl: PowerFlowView.cpp
\begin{lyxcode}
\_worker~=~std::thread(\&PowerFlowView::workerMethod,~this,~input);

\_progressThread~=~std::thread(\&PowerFlowView::progressMethod,~this);

gui::thread::asyncExecInMainThread({[}this,~displayData,~viewAlive{]}()

\{

~~~~if~(!{*}viewAlive)

~~~~~~~~return;

~~~~finishConversion(displayData);

\});
\end{lyxcode}
Ovaj isječak pokazuje osnovnu podjelu posla: konverzija ide u worker
thread, a progress se računa odvojeno. Završni prikaz rezultata se
prebacuje u glavni GUI thread, čime se poštuje natID pattern za siguran
GUI update.

\section{DLL plugin}

Ordinary DLL plugin predstavlja odvojeni C API sloj koji omogućava
pozivanje core konverzije bez direktnog korištenja GUI klase. Na taj
način se ista logika može koristiti iz glavne aplikacije i iz drugih
slojeva.

{\large\textcolor{red}{{[}OVDJE UBACITI SLIKU ILI DIJAGRAM: ordinary
DLL plugin tok{]}}}{\large\par}

\subsection{Kodni isječak: ordinary DLL API}

Fajl: TpfpcPluginApi.cpp
\begin{lyxcode}
int~tpfpcConvertThreePhase(const~TpfpcPluginInput{*}~input,~TpfpcPluginResult{*}~result)

\{

~~~~tpfpc::PowerFlowMatrixModel~matrixModel;

~~~~const~auto~conversionResult~=~matrixModel.convert(rectangularPower);

~~~~auto~polar~=~conversionResult.polarMatrix.getManipulator();

~~~~result->phaseA~=~makePluginPolarPower(polar(0,~0),~polar(0,~1));

~~~~return~1;

\}
\end{lyxcode}
Ovaj API prima ulaznu strukturu sa P/Q vrijednostima i puni izlaznu
strukturu sa polarnim rezultatima. Bitan je jer razdvaja core konverziju
od GUI-ja i omogućava plugin upotrebu.

\section{dTwin plugin}

dTwin import plugin implementira sc::IPlugin, eksportuje getPluginInterface
i u dTwin aplikaciji se pojavljuje pod Model -> Import kao TPFPC Default
Converter. Nakon unosa P/Q vrijednosti plugin generiše .dmodl model
koji sadrži rezultate konverzije.

Plugin se kopira u:
\begin{lyxcode}
C:
\end{lyxcode}
U toku lokalnog testiranja je utvrđeno da instalirani dTwin može učitati
i kopiju iz Program Files plugin foldera. Zbog toga je za demonstraciju
korištena ažurirana DLL kopija, uz prethodni backup stare DLL datoteke.
Ovo ne mijenja osnovni kod projekta.

{\large\textcolor{red}{{[}OVDJE UBACITI SLIKU: dTwin Model -> Import
meni{]}}}{\large\par}

{\large\textcolor{red}{{[}OVDJE UBACITI SLIKU: dTwin TPFPC plugin
prozor{]}}}{\large\par}

{\large\textcolor{red}{{[}OVDJE UBACITI SLIKU: generisani .dmodl model
u dTwin-u{]}}}{\large\par}

\subsection{Kodni isječak: dTwin export i callback}

Fajl: TpfpcDTwinPlugin.cpp
\begin{lyxcode}
extern~\textquotedbl C\textquotedbl{}

\{

PLUGIN\_API~sc::IPlugin{*}~getPluginInterface()

\{

~~~~return~\&s\_plugin;

\}

\}
\end{lyxcode}
Ovo je ulazna tačka koju dTwin loader traži u plugin DLL-u. Funkcija
vraća statički plugin objekat koji implementira sc::IPlugin interfejs.

\subsection{Kodni isječak: generisanje .dmodl i završetak import workflow-a}

Fajl: TpfpcDTwinPlugin.cpp
\begin{lyxcode}
const~auto~displayData~=~convertPowerFlow(rectangularPower);

const~auto~dmodlContent~=~createDmodlContent(displayData);

writeArchive(getArchive(ArchType::DigitalModel),~dmodlContent);

fileOut~<\textcompwordmark <~dmodlContent;

if~(\_onComplete)

~~~~\_onComplete(\_pPlugin);
\end{lyxcode}
Ovaj dio povezuje lokalnu konverziju sa dTwin import mehanizmom. .dmodl
sadržaj se zapisuje u digitalni archive i u izlazni fajl, a \_onComplete
obavještava dTwin da je import završen.

\section{Testiranje i validacija}

Za osnovnu provjeru koriste se balanced i unbalanced primjeri. Balanced
test provjerava da se faza A kopira u B i C, dok unbalanced test provjerava
da se sve tri faze mogu zadati nezavisno.

MATPOWER se koristi samo kao izvor referentnih P/Q podataka. Projekat
ne poredi rezultate sa cijelim MATPOWER solverom i nije MATPOWER parser.
Validirani dio je isključivo konverzija poznatih P/Q vrijednosti u
polarni oblik.

MATPOWER referentni primjer iz t\_case3p\_a:
\begin{lyxcode}
A:~P~=~1341.42,~Q~=~970.52,~mag~=~1655.6922,~angleDeg~=~35.8858

B:~P~=~2096.10,~Q~=~1341.41,~mag~=~2488.5771,~angleDeg~=~32.6174

C:~P~=~2672.34,~Q~=~1894.59,~mag~=~3275.8010,~angleDeg~=~35.3352
\end{lyxcode}
U projektu je dodat MatpowerReferenceTest. Test je pokrenut u Release
konfiguraciji i vratio je:
\begin{lyxcode}
MATPOWER~t\_case3p\_a~reference~conversion:~OK
\end{lyxcode}
{\large\textcolor{red}{{[}OVDJE UBACITI SLIKU ILI TABELU: MATPOWER
referentni test / tabela rezultata{]}}}{\large\par}

\section{Build i pokretanje}

Projekat se gradi preko CMake konfiguracije i Visual Studio/MSBuild
alata. Lokalno je potvrđen Windows Release ALL\_BUILD, uključujući
glavnu aplikaciju, ordinary DLL plugin, dTwin plugin, dTwin test executable
targete i MatpowerReferenceTest.

Release output folder:
\begin{lyxcode}
C:
\end{lyxcode}
Glavna aplikacija:
\begin{lyxcode}
SREES\_2026\_Bajramovic\_TPFPC.exe
\end{lyxcode}
dTwin plugin:
\begin{lyxcode}
SREES\_2026\_Bajramovic\_TPFPC\_dTwinPlugin.dll
\end{lyxcode}
Za pokretanje na drugom računaru potrebno je imati podešeno natID
okruženje ili pripremiti poseban Release paket sa potrebnim DLL zavisnostima.
Prema uslovima projekta, build treba dodatno potvrditi i na macOS-u
i na još jednom sistemu. Ta provjera nije potvrđena u ovom izvještaju
i navedena je kao planirano testiranje.

\section{Zaključak}

Projekat SREES\_2026\_Bajramovic\_TPFPC realizuje čisti trofazni P+jQ
u polarni konverter. Implementirani su core C++ proračun, natID GUI,
balanced i unbalanced režim, dense/sparse matrice, threading sa progress
indikatorom, ordinary DLL plugin, dTwin import plugin i MATPOWER referentna
validacija P/Q podataka.

Najvažnije ograničenje je da projekat nije puni load-flow solver i
ne pokušava zamijeniti MATPOWER ili implementirati Newton-Raphson
postupak. Njegova namjena je jasno određena: konverzija poznatih trofaznih
P/Q vrijednosti u polarne koordinate i prikaz tih rezultata kroz natID/dTwin
workflow.

\section{Literatura i reference}
\begin{enumerate}
\item natID SDK dokumentacija i lokalni primjeri.
\item dTwin plugin primjer iz lokalnog natID.Examples
\item MATPOWER three-phase dokumentacija: https://matpower.org/documentation/howto/three-phase.html
\item MATPOWER test cases repository: https://github.com/MATPOWER/matpower/tree/master/lib/t
\item CMake dokumentacija: https://cmake.org/documentation/
\item C++ standardna biblioteka, posebno funkcije sqrt i atan2 iz cmath.
\end{enumerate}

\end{document}
